% Section 5: Provisional Evidentiary Capabilities

\section{Provisional Evidentiary Capabilities}
\label{sec:primitives}

The diagnostics in Section~\ref{sec:diagnostics} generate five provisional
evidentiary capabilities. They are design hypotheses, not universal
requirements or a complete governance architecture. A capability may be
supplied by a vendor product, a bank-controlled system, a contractually required
export, an independent assessment, or a compensating process. The institution
remains responsible for claims it makes with vendor-supported evidence, but
responsibility does not imply unilateral power to alter the vendor's system.

Each capability is stated by the object it should expose and the limit on what
that object establishes. An examiner need not personally perform cryptographic
verification or model forensics: the review may inspect bank or third-party
validation, provenance, controls, and selected underlying evidence. Which
pathway is proportionate depends on the claim, risk, authority, and available
implementation.

\subsection{Capability 1: Revision-aware temporal capture}
\label{sec:primitives:temporal-capture}

The first capability preserves distinct, time-bound records rather than a
single mutable narrative. Its decision-time record captures the inputs and context
presented, observable actions, stated uncertainties, and provisional reasons.
Later records may add reflective interpretation or newly discovered reasons;
still later records may authorize, correct, or reassess the decision under
policy. Each layer carries its own timestamp, provenance, and binding to the
decision and prior record. Later reflection may contribute genuine
understanding, but it must not overwrite or impersonate the earlier account.

Tamper-evident commitments, controlled audit logs, signed records, or
independently maintained timestamps are possible implementations. A bank may
operate them or contract for them. An examiner can inspect validation of the
binding and selected records rather than reproduce the mechanism. Temporal
capture establishes when specified content was committed and whether revision
is detectable; it does not establish that the content was true, complete, or a
transparent account of human cognition.

\subsection{Capability 2: Evidence binding and exportability}
\label{sec:primitives:evidence-binding}

The second capability makes a declared relationship among a decision, artifact,
method, model and data versions, policy, review, and authorization traversable.
Inspection should be possible backward from a decision to its supporting
material and, where appropriate, forward from material to affected decisions.
The export must preserve identifiers, provenance, configurations, and relevant
validation rather than reduce the relationship to a product dashboard.

This is particularly important for community banks using vendor-controlled
systems. A bank must be able to obtain, directly or by contract, evidence
sufficient to support the claims it is obligated to make. A vendor may provide
a structured export; a bank may maintain external bindings; an independent
assessor may validate proprietary components; or a compensating control may
narrow the claim. Product access without inspectable evidence leaves the
coordination gap in place.

Binding an explanation to a decision establishes association and provenance.
It does not prove causal influence, method fidelity, or evidentiary sufficiency.
Those are separate propositions tested with method-appropriate evidence.

\subsection{Capability 3: Non-collapsing uncertainty representation}
\label{sec:primitives:confidence-characterization}

The third capability preserves materially different reasons for uncertainty
instead of forcing them into one probability or confidence label. Depending on
the decision, a useful representation may distinguish supporting evidence,
counterevidence, conflict among sources, model or measurement uncertainty,
missing information, and ambiguity about applicable policy. A single value can
be well calibrated for one target while remaining silent about which of these
conditions produced it.

The paper does not prescribe a fixed number of dimensions or one mathematical
formalism. A structured record, interval, set-valued output, qualitative state,
or combination may be appropriate. The functional test is whether material
support, counterevidence, conflict, and insufficiency survive reporting and can
be bound to the resulting action. Vendors may expose native uncertainty and
provenance; banks must decide how policy treats the represented state; and
validation should test whether the representation is stable, meaningful, and
used as declared. The representation does not itself establish that the
underlying evidence is accurate.

\subsection{Capability 4: Technical-change monitoring}
\label{sec:primitives:technical-change}

Technical-change monitoring asks whether the empirical system has changed.
Relevant objects include input distributions, relationships between inputs and
targets, measured performance, calibration, subgroup behavior, and the model or
data pipeline itself. Data-drift tests, concept-drift analysis, outcome
monitoring, version comparison, and targeted validation answer different
questions; no shared detector is presumed.

The likely implementation path combines vendor telemetry and change notices
with bank or third-party validation against the institution's operational
population. Examiners can inspect the monitoring design, thresholds, escalation
records, and selected results. A detected change directs investigation; it does
not establish cause, harm, or noncompliance.

\subsection{Capability 5: Institutional-conformance monitoring}
\label{sec:primitives:institutional-conformance}

Institutional-conformance monitoring asks a different question: whether
operational practice continues to match declared policy, authority, and control
design. Its evidence includes workflow and override records, access and
authorization histories, exceptions, complaints, sampled case files,
attestation-scope reconciliation, and changes in policy interpretation.
Statistical drift may be absent while practice diverges from policy, or present
while practice remains conformant.

The bank ordinarily owns the policy relationship, while vendors must expose the
events and configurations needed to test it. Internal audit, compliance,
independent review, affected-party evidence, and examiner-selected samples can
challenge a self-selected institutional account. A discrepancy establishes a
conformance question and may support a finding where applicable authority and
evidence do so; the monitoring capability alone does not determine legal
significance.

\subsection{What the capabilities contribute}
\label{sec:primitives:compound}

Together, the five capabilities can make the diagnostics answerable:
revision-aware capture preserves temporal layers; binding and exportability make
declared relations inspectable; non-collapsing representation preserves the
kind of uncertainty at issue; and the two monitoring capabilities distinguish
technical change from institutional divergence. Multiple coherent
implementations may exist, and a compensating process may sometimes provide the
required evidence more proportionately than new infrastructure.

Their presence does not establish that every affected interest is represented,
and their absence does not prove that representation is impossible. Cost,
security, confidentiality, records law, and vendor leverage constrain the
design space. The practical question is whether a proposed implementation makes
the relevant artifact-to-inference relation more inspectable at a proportionate
cost. That question remains empirical.
